\section{Private-Trajectory Methodology}
\label{app:private-trace}

The private case is one self-hosted longitudinal lineage, not a sample of
independent tasks.  A fixed cutoff selects one root and every recursive spawn
descendant created before that cutoff.  For the two observed Codex versions,
pinned source says the legacy fork filter removes inter-agent metadata from
copied history~\cite{codexspawn145,codexspawn146}.  The extractor identifies
each child's native suffix from that source-grounded delivery boundary, reads
one immutable byte image per rollout, checks header, edge, and cutoff
consistency, and fails closed on ambiguous lineage or unsupported versions.
Two unchanged-input extractions were byte-identical.  Synthetic regressions
cover cutoff placement, false spawn triggers, multiple parents, external-root
parents, malformed rows, duplicate tool IDs, determinism, and the allowlisted
privacy schema.

One record carries a UI rollback marker, but it is not treated as a typed
Restore; copied rows are not duplicate effects, and tool-call IDs are not
durable redemption receipts.  Claude contributes only an explicitly selected
historical project index.  Its declared raw action paths were unavailable;
bounded exact-ID recovery found metadata references but no provenance-valid
action corpus, so we report no Claude tool, mutation, fork, or effect counts.

The author-operated case recruited or manipulated no users.  Extraction emits
only path-free allowlisted aggregates and private keyed commitments; it excludes
prompts, commands, bodies, paths, and stable raw identifiers.  Raw histories,
the commitment key, and the private summary are withheld.  The published
artifact contains the extractor, public-data analysis, formal development,
compiler/verifier, and synthetic fixtures.  These choices reduce linkage risk
while retaining the limited workload-shape evidence.
